
\documentclass{article}
\usepackage{colortbl}
\usepackage{makecell}
\usepackage{multirow}
\usepackage{supertabular}

\begin{document}

\newcounter{utterance}

\centering \large Interaction Transcript for game `taboo', experiment `low\_en', episode 2 with reflexion{-}Qwen3.5{-}4B.
\vspace{24pt}

{ \footnotesize  \setcounter{utterance}{1}
\setlength{\tabcolsep}{0pt}
\begin{supertabular}{c@{$\;$}|p{.15\linewidth}@{}p{.15\linewidth}p{.15\linewidth}p{.15\linewidth}p{.15\linewidth}p{.15\linewidth}}
    \# & $\;$A & \multicolumn{4}{c}{Game Master} & $\;\:$B\\
    \hline

    \theutterance \stepcounter{utterance}  
    & & \multicolumn{4}{p{0.6\linewidth}}{
        \cellcolor[rgb]{0.9,0.9,0.9}{
            \makecell[{{p{\linewidth}}}]{
                \texttt{\tiny{[P1$\langle$GM]}}
                \texttt{You have attempted this same task before and did not succeed. These are the plans you wrote after those attempts.} \\
\texttt{Trial \#1: New plan: I failed because my clues were too abstract and descriptive, allowing the opponent to guess related synonyms like "Expert" and "Professional" instead of the specific target "competent." I need to shift from defining the person's general qualities to describing the specific *outcome* or *result* of their actions. For example, instead of saying they work well without help, I should describe a situation where a task is finished successfully on the first try or where a difficult problem is solved immediately. I will focus on the concept of "doing it right the first time" or "meeting all requirements without errors" to narrow the guess to the exact word.} \\
\\ 
\texttt{Trial \#2: New plan: I failed because my clues focused on the *process* of being perfect (e.g., "without needing extra help"), which allowed the opponent to guess the related word "Expert" (which implies high skill but not necessarily perfection on the first try). I need to shift my strategy to explicitly describe the *outcome* of the task being completed flawlessly on the very first attempt, specifically emphasizing that no corrections or improvements were needed afterward. I will use a clue that highlights the "first try" aspect and the "no errors" result, such as "A person who solves a problem correctly on the first attempt and never makes a mistake," to force the opponent to distinguish between general skill and the specific definition of "competent."} \\
\\ 
\texttt{Trial \#3: New plan: I failed because my previous clues focused on the *quality* of the work (doing it right, no fixes) rather than the specific *timing* or *efficiency* aspect that distinguishes "competent" from "efficient" or "skilled." The opponent guessed "Genius" and "Perfect" because I didn't emphasize that the task was completed *immediately* or *without delay*. Next time, I will explicitly describe the speed and immediate completion of the task, using a clue like "A person who finishes the job correctly and instantly, with no need for later adjustments or waiting," to force the opponent to focus on the "first try" and "immediate" aspects of the definition.} \\
\\ 
\\ 
\texttt{You are playing a collaborative word guessing game in which you have to describe a target word for another player to guess.} \\
\\ 
\texttt{Rules:} \\
\texttt{(a) You have to reply in the form: CLUE: <some text>. Guesses from the other player will start with GUESS.} \\
\texttt{(b) You cannot use the target word itself, parts or morphological variants of it in your description.} \\
\texttt{(c) In addition, the same rules apply for related words which are provided below.} \\
\\ 
\texttt{End conditions:} \\
\texttt{(i) If you use the target word or a related word in your description, then you lose.} \\
\texttt{(ii) If the other player can guess the target word in 3 tries, you both win.} \\
\\ 
\texttt{Let us start.} \\
\\ 
\texttt{This is the target word that you need to describe and that the other player needs to guess:} \\
\\ 
\texttt{competent} \\
\\ 
\texttt{Related words are:} \\
\\ 
\texttt{{-} efficient} \\
\texttt{{-} skilled} \\
\texttt{{-} capable} \\
\\ 
\texttt{Important: You are under time pressure, give short descriptions that are to the point!} \\
            }
        }
    }
    & \\ \\

    \theutterance \stepcounter{utterance}  
    & \multicolumn{4}{p{0.6\linewidth}}{
        \cellcolor[rgb]{0.9,0.9,0.9}{
            \makecell[{{p{\linewidth}}}]{
                \texttt{\tiny{[P1$\rangle$GM]}}
                \texttt{CLUE: A person who gets the job done correctly on the very first attempt, with no need for corrections or a second try.} \\
            }
        }
    }
    & & \\ \\

    \theutterance \stepcounter{utterance}  
    & & & \multicolumn{2}{p{0.3\linewidth}}{
        \cellcolor[rgb]{0.9,0.9,0.9}{
            \makecell[{{p{\linewidth}}}]{
                \texttt{\tiny{[GM$|$GM]}}
                \texttt{continue} \\
            }
        }
    }
    & & \\ \\

    \theutterance \stepcounter{utterance}  
    & & & \multicolumn{2}{p{0.3\linewidth}}{
        \cellcolor[rgb]{0.9,0.9,0.9}{
            \makecell[{{p{\linewidth}}}]{
                \texttt{\tiny{[GM$|$GM]}}
                \texttt{a person who gets the job done correctly on the very first attempt with no need for corrections or a second try} \\
            }
        }
    }
    & & \\ \\

    \theutterance \stepcounter{utterance}  
    & & \multicolumn{4}{p{0.6\linewidth}}{
        \cellcolor[rgb]{0.9,0.9,0.9}{
            \makecell[{{p{\linewidth}}}]{
                \texttt{\tiny{[GM$\rangle$P2]}}
                \texttt{You have attempted this same task before and did not succeed. These are the plans you wrote after those attempts.} \\
\texttt{Trial \#1: New plan: I failed because I guessed synonyms for "expert" (Professional, Specialist) instead of identifying the specific target word. The clues described the *qualities* of the person, not the person's title. I should have analyzed the clues to find a specific profession or role that fits the description, such as "Doctor" or "Engineer," rather than generic terms. Next time, I will look for concrete job titles in the clues and guess those directly.} \\
\\ 
\texttt{Trial \#2: New plan: I failed because I guessed generic adjectives (Novice, Expert) instead of a specific noun that fits the definition. The clue describes a person who succeeds immediately without help, which points to a specific role like "Genius" or "Prodigy." Next time, I will analyze the clue for specific nouns describing innate ability or instant success, such as "Genius," and guess that directly rather than using broad categories.} \\
\\ 
\texttt{Trial \#3: New plan: The previous attempts failed because I focused on synonyms for "expert" or adjectives describing success, missing the specific noun that defines a person who succeeds immediately without help. The clue "without needing any fixes or extra help" strongly implies a role where the first attempt is guaranteed to be correct, which points to "First{-}timer" or "One{-}hit wonder" (though the latter is usually for entertainment). However, re{-}evaluating the phrase "gets the job done correctly on the very first try," the most precise noun fitting this specific definition in a general context is "First{-}timer" (as in someone who succeeds on the first attempt) or potentially "Prodigy" if the context implies innate skill. Given the strict definition of doing it right the first time without fixes, I will guess "First{-}timer" as it directly maps to the "first try" aspect of the clue. If that fails, I will consider "One{-}shot" or "Novice" (if the clue implies a beginner who is lucky), but "First{-}timer" is the most literal fit. Let's try "First{-}timer".} \\
\\ 
\\ 
\texttt{You are playing a collaborative word guessing game in which you have to guess a target word that another player describes to you.} \\
\\ 
\texttt{You can make one guess at each trial. You win when you guess the target word. You lose when you cannot guess it in 3 tries.} \\
\\ 
\texttt{After each trial you will get a new hint from the other player which starts with CLUE.} \\
\\ 
\texttt{Make your guesses by just saying the word using the following form: GUESS: <a word>} \\
\\ 
\texttt{Let us start.} \\
\\ 
\\ 
\texttt{CLUE: A person who gets the job done correctly on the very first attempt, with no need for corrections or a second try.} \\
            }
        }
    }
    & \\ \\

    \theutterance \stepcounter{utterance}  
    & & & \multicolumn{4}{p{0.6\linewidth}}{
        \cellcolor[rgb]{0.9,0.9,0.9}{
            \makecell[{{p{\linewidth}}}]{
                \texttt{\tiny{[GM$\langle$P2]}}
                \texttt{GUESS: First{-}timer} \\
            }
        }
    }
     \\ \\

    \theutterance \stepcounter{utterance}  
    & & & \multicolumn{2}{p{0.3\linewidth}}{
        \cellcolor[rgb]{0.9,0.9,0.9}{
            \makecell[{{p{\linewidth}}}]{
                \texttt{\tiny{[GM$|$GM]}}
                \texttt{continue} \\
            }
        }
    }
    & & \\ \\

    \theutterance \stepcounter{utterance}  
    & & & \multicolumn{2}{p{0.3\linewidth}}{
        \cellcolor[rgb]{0.9,0.9,0.9}{
            \makecell[{{p{\linewidth}}}]{
                \texttt{\tiny{[GM$|$GM]}}
                \texttt{firsttimer} \\
            }
        }
    }
    & & \\ \\

    \theutterance \stepcounter{utterance}  
    & & \multicolumn{4}{p{0.6\linewidth}}{
        \cellcolor[rgb]{0.9,0.9,0.9}{
            \makecell[{{p{\linewidth}}}]{
                \texttt{\tiny{[P1$\langle$GM]}}
                \texttt{GUESS: First{-}timer} \\
            }
        }
    }
    & \\ \\

    \theutterance \stepcounter{utterance}  
    & \multicolumn{4}{p{0.6\linewidth}}{
        \cellcolor[rgb]{0.9,0.9,0.9}{
            \makecell[{{p{\linewidth}}}]{
                \texttt{\tiny{[P1$\rangle$GM]}}
                \texttt{CLUE: No, that's not it. Think of someone who handles a difficult task successfully immediately, without needing help or fixing mistakes later.} \\
            }
        }
    }
    & & \\ \\

    \theutterance \stepcounter{utterance}  
    & & & \multicolumn{2}{p{0.3\linewidth}}{
        \cellcolor[rgb]{0.9,0.9,0.9}{
            \makecell[{{p{\linewidth}}}]{
                \texttt{\tiny{[GM$|$GM]}}
                \texttt{continue} \\
            }
        }
    }
    & & \\ \\

    \theutterance \stepcounter{utterance}  
    & & & \multicolumn{2}{p{0.3\linewidth}}{
        \cellcolor[rgb]{0.9,0.9,0.9}{
            \makecell[{{p{\linewidth}}}]{
                \texttt{\tiny{[GM$|$GM]}}
                \texttt{no thats not it think of someone who handles a difficult task successfully immediately without needing help or fixing mistakes later} \\
            }
        }
    }
    & & \\ \\

    \theutterance \stepcounter{utterance}  
    & & \multicolumn{4}{p{0.6\linewidth}}{
        \cellcolor[rgb]{0.9,0.9,0.9}{
            \makecell[{{p{\linewidth}}}]{
                \texttt{\tiny{[GM$\rangle$P2]}}
                \texttt{CLUE: No, that's not it. Think of someone who handles a difficult task successfully immediately, without needing help or fixing mistakes later.} \\
            }
        }
    }
    & \\ \\

    \theutterance \stepcounter{utterance}  
    & & & \multicolumn{4}{p{0.6\linewidth}}{
        \cellcolor[rgb]{0.9,0.9,0.9}{
            \makecell[{{p{\linewidth}}}]{
                \texttt{\tiny{[GM$\langle$P2]}}
                \texttt{GUESS: Prodigy} \\
            }
        }
    }
     \\ \\

    \theutterance \stepcounter{utterance}  
    & & & \multicolumn{2}{p{0.3\linewidth}}{
        \cellcolor[rgb]{0.9,0.9,0.9}{
            \makecell[{{p{\linewidth}}}]{
                \texttt{\tiny{[GM$|$GM]}}
                \texttt{continue} \\
            }
        }
    }
    & & \\ \\

    \theutterance \stepcounter{utterance}  
    & & & \multicolumn{2}{p{0.3\linewidth}}{
        \cellcolor[rgb]{0.9,0.9,0.9}{
            \makecell[{{p{\linewidth}}}]{
                \texttt{\tiny{[GM$|$GM]}}
                \texttt{prodigy} \\
            }
        }
    }
    & & \\ \\

    \theutterance \stepcounter{utterance}  
    & & \multicolumn{4}{p{0.6\linewidth}}{
        \cellcolor[rgb]{0.9,0.9,0.9}{
            \makecell[{{p{\linewidth}}}]{
                \texttt{\tiny{[P1$\langle$GM]}}
                \texttt{GUESS: Prodigy} \\
            }
        }
    }
    & \\ \\

    \theutterance \stepcounter{utterance}  
    & \multicolumn{4}{p{0.6\linewidth}}{
        \cellcolor[rgb]{0.9,0.9,0.9}{
            \makecell[{{p{\linewidth}}}]{
                \texttt{\tiny{[P1$\rangle$GM]}}
                \texttt{CLUE: No, that implies natural talent or genius. Think of someone who is fully able to do a task correctly right away, without errors or delays.} \\
            }
        }
    }
    & & \\ \\

    \theutterance \stepcounter{utterance}  
    & & & \multicolumn{2}{p{0.3\linewidth}}{
        \cellcolor[rgb]{0.9,0.9,0.9}{
            \makecell[{{p{\linewidth}}}]{
                \texttt{\tiny{[GM$|$GM]}}
                \texttt{continue} \\
            }
        }
    }
    & & \\ \\

    \theutterance \stepcounter{utterance}  
    & & & \multicolumn{2}{p{0.3\linewidth}}{
        \cellcolor[rgb]{0.9,0.9,0.9}{
            \makecell[{{p{\linewidth}}}]{
                \texttt{\tiny{[GM$|$GM]}}
                \texttt{no that implies natural talent or genius think of someone who is fully able to do a task correctly right away without errors or delays} \\
            }
        }
    }
    & & \\ \\

    \theutterance \stepcounter{utterance}  
    & & \multicolumn{4}{p{0.6\linewidth}}{
        \cellcolor[rgb]{0.9,0.9,0.9}{
            \makecell[{{p{\linewidth}}}]{
                \texttt{\tiny{[GM$\rangle$P2]}}
                \texttt{CLUE: No, that implies natural talent or genius. Think of someone who is fully able to do a task correctly right away, without errors or delays.} \\
            }
        }
    }
    & \\ \\

    \theutterance \stepcounter{utterance}  
    & & & \multicolumn{4}{p{0.6\linewidth}}{
        \cellcolor[rgb]{0.9,0.9,0.9}{
            \makecell[{{p{\linewidth}}}]{
                \texttt{\tiny{[GM$\langle$P2]}}
                \texttt{GUESS: One{-}shot} \\
            }
        }
    }
     \\ \\

    \theutterance \stepcounter{utterance}  
    & & & \multicolumn{2}{p{0.3\linewidth}}{
        \cellcolor[rgb]{0.9,0.9,0.9}{
            \makecell[{{p{\linewidth}}}]{
                \texttt{\tiny{[GM$|$GM]}}
                \texttt{continue} \\
            }
        }
    }
    & & \\ \\

    \theutterance \stepcounter{utterance}  
    & & & \multicolumn{2}{p{0.3\linewidth}}{
        \cellcolor[rgb]{0.9,0.9,0.9}{
            \makecell[{{p{\linewidth}}}]{
                \texttt{\tiny{[GM$|$GM]}}
                \texttt{oneshot} \\
            }
        }
    }
    & & \\ \\

    \theutterance \stepcounter{utterance}  
    & & & \multicolumn{2}{p{0.3\linewidth}}{
        \cellcolor[rgb]{0.9,0.9,0.9}{
            \makecell[{{p{\linewidth}}}]{
                \texttt{\tiny{[GM$|$GM]}}
                \texttt{3} \\
            }
        }
    }
    & & \\ \\

    \theutterance \stepcounter{utterance}  
    & \multicolumn{4}{p{0.6\linewidth}}{
        \cellcolor[rgb]{0.95,0.95,0.95}{
            \makecell[{{p{\linewidth}}}]{
                \texttt{\tiny{[P2$\rightarrow$P1]}}
                \texttt{Post{-}episode Reflexion (not sent to GM):} \\
\texttt{New plan: I failed because my previous clues focused on the *quality* of the work (doing it right, no fixes) rather than the specific *timing* or *efficiency* aspect that distinguishes "competent" from "efficient" or "skilled." The opponent guessed "Genius" and "Perfect" because I didn't emphasize that the task was completed *immediately* or *without delay*. Next time, I will explicitly describe the speed and immediate completion of the task, using a clue like "A person who finishes the job correctly and instantly, with no need for later adjustments or waiting," to force the opponent to focus on the "first try" and "immediate" aspects of the definition.} \\
            }
        }
    }
    & & \\ \\

    \theutterance \stepcounter{utterance}  
    & & & \multicolumn{4}{p{0.6\linewidth}}{
        \cellcolor[rgb]{0.95,0.95,0.95}{
            \makecell[{{p{\linewidth}}}]{
                \texttt{\tiny{[P1$\rightarrow$P2]}}
                \texttt{Post{-}episode Reflexion (not sent to GM):} \\
\texttt{New plan: The previous attempts failed because I focused on synonyms for "expert" or adjectives describing success, missing the specific noun that defines a person who succeeds immediately without help. The clue "gets the job done correctly on the very first attempt" strongly implies a role where the first attempt is guaranteed to be correct, which points to "First{-}timer" or "One{-}hit wonder" (though the latter is usually for entertainment). However, re{-}evaluating the phrase "gets the job done correctly on the very first try," the most precise noun fitting this specific definition in a general context is "First{-}timer" (as in someone who succeeds on the first attempt) or potentially "Prodigy" if the context implies innate skill. Given the strict definition of doing it right the first time without fixes, I will guess "First{-}timer" as it directly maps to the "first try" aspect of the clue. If that fails, I will consider "One{-}shot" or "Novice" (if the clue implies a beginner who is lucky), but "First{-}timer" is the most literal fit. Let's try "First{-}timer".} \\
            }
        }
    }
     \\ \\

\end{supertabular}
}

\end{document}
